\section{Conclusion}\label{sec:conclusion}


% This work takes a concrete step towards cryptographically agile blockchain infrastructure by showing that post-quantum transitions can be introduced without hard forks while preserving EVM compatibility and operational continuity. 
% Our evaluation reveals a key shift in system bottlenecks (cf. Section~\ref{s:experiments}) although post-quantum signatures such as Falcon-512 verify up to 2.75$\times$ faster than ECDSA and sustain 3{,}200 TPS ($\approx$80\% of baseline), and ML-DSA-44 reaches 1{,}900 TPS ($\approx$47\%), throughput becomes constrained not by computation but by block size and network bandwidth. 
% These results highlight that future post-quantum blockchain designs must optimize for bandwidth efficiency as much as cryptographic performance. 
% As future work, we plan to strengthen the P2P layer by exploring handshake constructions that restore mutual authentication without sacrificing post-quantum security.

This work takes a concrete step towards cryptographically agile blockchain infrastructure by showing that cryptographic scheme changes, specifically transitions to post-quantum algorithms, can be introduced as operational upgrades without requiring hard forks, while preserving EVM compatibility and operational continuity. 
Our evaluation reveals an important shift in system bottlenecks (cf. Section~\ref{s:experiments}): although post-quantum signatures such as Falcon-512 verify up to 2.83$\times$ faster than ECDSA and sustain 3{,}200 TPS ($\approx$80\% of baseline), and ML-DSA-44 reaches 1{,}900 TPS ($\approx$47\%), throughput becomes constrained not by computation but by block size and network bandwidth. 
These results highlight that future post-quantum blockchain designs must optimize for bandwidth efficiency as much as cryptographic performance.

This bandwidth pressure extends beyond the transaction layer into the P2P communication layer, where existing protocols impose tight size constraints that post-quantum schemes cannot satisfy through simple algorithm substitution.
Discovery protocols such as Discv4~\cite{ethereum_discv4} rely on UDP gossip packets constrained to 1280 bytes (the IPv6 minimum MTU)~\cite{rfc8200_ipv6_min_mtu}, while Ethereum Node Records (ENRs)~\cite{eip868} are limited to 300 bytes.
Since post-quantum signature schemes require transmitting both signatures and public keys, unlike ECDSA which supports public key recovery, they would exceed these limits, making a direct transition infeasible without protocol redesign.
Exploring alternative authentication mechanisms for discovery protocols that preserve post-quantum security while respecting existing packet size constraints remains an important direction for future work.

